\section{Classification of the four-dimensional static base}
\label{q4:cla:section}

We now classify the static base obtained in
Proposition~\ref{q4:sta:complete-base}.  The positive set $\mathcal U=\{N>0\}$
is connected and contains the boundary collar.  On $\mathcal U$ the metric
$h_b$ and the lapse $\lambda=\sqrt N$ satisfy
\begin{equation}\label{q4:cla:static-equations}
 \lambda\Ric_{h_b}=\Hess_{h_b}\lambda,
 \qquad \Delta_{h_b}\lambda=0,
 \qquad \Scal_{h_b}=0,
 \qquad 0<\lambda<1.
\end{equation}
The base is complete with $S$ attached; approaches to interior zeros of $N$
have infinite base length.  Its boundary data are
\begin{equation}\label{q4:cla:boundary-data}
 h_b|_{TS}=g|_{TS},\qquad
 \lambda|_S=0,\qquad
 \partial_{\nu_{h_b}}\lambda=\kappa,\qquad
 \mathrm{II}_{h_b}=0,
 \qquad \kappa=\frac1{\sqrt{2m}}.
\end{equation}
These are conclusions of the preceding argument, not hypotheses on the
original data.  In particular, neither simple connectivity nor absence
of additional complete ends of the base is available yet.

Our goal is both the explicit static metric and the identity
$\mathcal U=\operatorname{int}\Omega$ in
Proposition~\ref{q4:cla:static-classification}.  The argument allows additional
complete ends until Euclidean rigidity of the conformal double excludes
interior null regions and identifies the whole exterior.

\subsection{Improvement of the static asymptotics}

After a constant linear change of the coordinates at infinity, the
positive limiting matrix of $h_b$ becomes the identity.  For a fixed
$q_0\in(1,\min\{q,2\})$ we then have
$h_b-\delta=O_2(r^{-q_0})$ and
$\lambda-1=O_2(r^{-q_0})$.  Introduce
\begin{equation}\label{q4:cla:reduced-variables}
 U=\log\lambda,\qquad \gamma=\lambda h_b.
\end{equation}
The conformal length factor from $h_b$ to $\gamma$ is $e^{U/2}$.
The four-dimensional Ricci transformation and
\eqref{q4:cla:static-equations} give
\begin{equation}\label{q4:cla:reduced-equations}
 \Ric_\gamma=\frac32\,\dd U\otimes\dd U,
 \qquad \Delta_\gamma U=0.
\end{equation}
Indeed,
$\Ric_{h_b}=\Hess_{h_b}U+\dd U\otimes\dd U$ and
$\Delta_{h_b}U=-|\dd U|_{h_b}^2$, which yield both identities directly.

\begin{lemma}[Static end expansion]\label{q4:cla:asymptotics}
There are coordinates harmonic for $\gamma$ on a sufficiently distant
part of the end, a number $\epsilon\in(0,1)$, a constant $M_1>0$, and
homogeneous harmonic polynomials $H_1,H_2$ of degrees one and two such
that, for every fixed nonnegative integer $k$,
\begin{align}
 \gamma-\delta&=O_k(r^{-2-\epsilon}),\label{q4:cla:gamma-decay}\\
 U&=-\frac{M_1}{r^2}
       +\frac{H_1(x)}{r^4}
       +\frac{H_2(x)}{r^6}
       +O_k(r^{-4-\epsilon}).\label{q4:cla:U-expansion}
\end{align}
In addition $\gamma-\delta=O_k(r^{-3})$.
\end{lemma}

\begin{proof}
We give the coordinate and expansion arguments, since the initial
falloff alone would not justify the compactification below.

Extend $\gamma$ from a sufficiently distant region to a smooth metric
on $\R^4$, using a cutoff interpolation with $\delta$.  Write its
Laplacian as
\[
 \Delta_\gamma
 =\Delta_\delta+D^{ij}\partial_i\partial_j+D^i\partial_i.
\]
By making the interpolation sufficiently far out, the scaled
$C^{0,\alpha}$ norms of $D^{ij}$ and $(1+r)D^i$ are as small as desired,
for any fixed $\alpha\in(0,1)$.  Here a scaled norm means the ordinary
norm after rescaling each annulus to fixed size.  The initial
differentiated falloff gives the needed Hölder estimates: second
derivatives bound the scaled Hölder seminorm of first derivatives.
Moreover $D^i=O(r^{-1-q_0})$ in these scaled norms.

Put $a=1-q_0\in(-1,0)$.  On all of $\R^4$, with weight $1+r$, the
Euclidean Newton operator maps scaled $C^{0,\alpha}$ functions of order
$a-2$ to scaled $C^{2,\alpha}$ functions of order $a$.
For completeness, the zeroth-order estimate follows from the kernel
$C|x-y|^{-2}$ by splitting the integral into $|y|<r/2$, a comparable
annulus, and $|y|>2r$; convergence uses $-2<a<0$.  Interior elliptic
estimates after rescaling give the two derivative and Hölder bounds.
The small coefficient norms therefore make the map obtained from
\[
 \Delta_\delta v^i
 =-D^i-D^{ab}\partial_a\partial_b v^i-D^a\partial_a v^i
\]
a contraction in that weighted space.  Its solution satisfies
$v=O(r^{1-q_0})$ with the stated scaled $C^{2,\alpha}$ bounds, and
$x^i+v^i$ are $\gamma$-harmonic coordinates on the far end.
They are indeed coordinates there: their first derivative tends to
the identity, and a cutoff of the correction can be chosen to have
globally small first derivative.  Local elliptic regularity makes the
coordinates smooth wherever the original metric is smooth.

In the harmonic chart, $\gamma-\delta$ and $U$ initially have order
$-q_0$ in scaled $C^{1,\alpha}$: one derivative of the transformed
metric uses two derivatives of the coordinate correction.  We do not
assume a weighted second-derivative estimate at this stage.  The chart
and the fields are nevertheless smooth locally, so
Equations~\eqref{q4:cla:reduced-equations}
take the coupled elliptic form
\begin{equation}\label{q4:cla:harmonic-system}
 -\frac12\gamma^{ab}\partial_a\partial_b\gamma_{ij}
       +Q_{ij}(\gamma,\partial\gamma)
       =\frac32 U_iU_j,
 \qquad \gamma^{ab}\partial_a\partial_bU=0,
\end{equation}
where $Q$ is quadratic in first derivatives, with smooth coefficients
depending on $\gamma$.  On an annulus rescaled to unit size, the initial
$C^{1,\alpha}$ norms of $\gamma-\delta$ and $U$ are $O(R^{-q_0})$.
Both quadratic right sides have $C^{0,\alpha}$ norm $O(R^{-2q_0})$,
and the principal coefficients are uniformly elliptic with bounded
$C^{1,\alpha}$ norm.  Interior Schauder estimates on a smaller annulus
therefore restore $C^{2,\alpha}$ bounds of order $R^{-q_0}$ for both
unknowns.  Differentiating this coupled system now inductively gives
$\gamma-\delta,U=O_k(r^{-q_0})$ for every fixed $k$.
Thus no higher derivative falloff is being assumed in the original
coordinates.  It follows in particular that
\begin{equation}\label{q4:cla:first-euclidean-source}
 \Delta_\delta(\gamma-\delta),\ \Delta_\delta U
       =O_k(r^{-2-2q_0}).
\end{equation}

We use the following elementary multipole fact in dimension four.
If a decaying smooth function $w$ on an end satisfies
\[
 \Delta_\delta w=O_k(r^{-4-j-\eta}),
 \qquad j\in\{0,1,2,\ldots\},\quad 0<\eta<1,
\]
then
\begin{equation}\label{q4:cla:multipole}
 w=\sum_{d=0}^j\frac{P_d(x)}{r^{2+2d}}
       +O_k(r^{-2-j-\eta}),
\end{equation}
with $P_d$ homogeneous harmonic of degree $d$.  To prove this, cut $w$
off to the end and represent the result by the whole-space Newton
potential of its Laplacian.  The difference is an entire harmonic
function tending to zero, hence vanishes.  The source has convergent
moments through degree $j$.  Taylor expansion of the kernel in the
region $|y|<r/2$ has remainder bounded by
$C r^{-3-j}|y|^{j+1}$; integration, and estimation of the complementary
region after subtraction of the same moments, gives the error in
\eqref{q4:cla:multipole}.  The local singularity of the Newton kernel is
integrable.  Rescaled elliptic estimates supply the differentiated
error bounds.  This also proves the statement simultaneously for
each fixed finite number of derivatives.

Choose $\epsilon\in(0,1)$ with $4+\epsilon<2+2q_0$.
Applying \eqref{q4:cla:multipole} first with $j=0$ gives monopole
expansions for $\gamma-\delta$ and $U$.  Write the metric monopole as
$D_{ij}/r^2$.  The harmonic-coordinate condition is
\[
 \partial_j\bigl(\sqrt{\det\gamma}\,\gamma^{ji}\bigr)=0.
\]
Its leading homogeneous term implies
\[
 \bigl(D-\tfrac12(\tr D)\Id\bigr)x=0
       \quad\hbox{for every }x.
\]
Taking the trace in dimension four gives $D=0$.
This proves \eqref{q4:cla:gamma-decay}.  Since $U=O_k(r^{-2})$, its
equation now improves to
$\Delta_\delta U=O_k(r^{-6-\epsilon})$.
The case $j=2$ of \eqref{q4:cla:multipole} gives
\eqref{q4:cla:U-expansion}, with the sign of its constant temporarily
undetermined.  In particular, the strict exponent in this source
estimate excludes a logarithmic term at quadrupole order.  The first
use of \eqref{q4:cla:multipole} likewise accounts for the entire order
$r^{-2}$ tensor term, rather than only its spherical average.
The metric equation in
\eqref{q4:cla:harmonic-system} has Euclidean source $O_k(r^{-6})$ at this
stage.  The case $j=1$ of \eqref{q4:cla:multipole}, with any smaller
positive exponent if necessary, gives
$\gamma-\delta=O_k(r^{-3})$.
No expansion of the metric beyond this order is asserted or needed.

Finally $-U$ is a positive $\gamma$-harmonic function.
On a sufficiently distant region the function
$r^{-2}+r^{-2-\epsilon}$ is positive and subharmonic for $\gamma$;
its favorable Euclidean Laplacian dominates the metric error.
A small positive multiple lies below $-U$ on a fixed large sphere,
and both tend to zero at infinity.  The maximum principle therefore
keeps it below $-U$ throughout the region.  Taking the limit of
$r^2(-U)$ proves $M_1>0$.
\end{proof}

\begin{remark}\label{q4:cla:base-mass}
The coefficient in Lemma~\ref{q4:cla:asymptotics} has the expected
four-dimensional mass normalization:
\[
 h_b=e^{-U}\gamma
       =(1+M_1/r^2)\delta+O_k(r^{-3}),\qquad
 \lambda=1-M_1/r^2+O_k(r^{-3}).
\]
Direct substitution in the energy flux with factor $1/(6\omega_3)$
gives base energy $M_1$.  We have not identified that chart with the
original ADM frame.  The equality $M_1=m$ will instead follow from
the boundary constant $\kappa$.
\end{remark}

The static equations have improved the initial weak decay enough for
two different operations: the plus conformal metric will have zero mass,
and the minus metric will extend through its point at infinity with two
continuous derivatives. The dipole and quadrupole terms are retained
because a leading mass expansion alone would not justify the latter.

\subsection{The complete conformal double}

We use the conformal-doubling strategy of Bunting and Masood-ul-Alam
\cite{BuntingMasoodUlAlam1987} and its higher-dimensional form due
to Gibbons, Ida, and Shiromizu
\cite[Section~2]{GibbonsIdaShiromizu2002}. The static asymptotics
retain a zero-mass end on one copy and compactify that end on the
other. We keep all possible complete ends arising from interior
zeros of $N$ until the separate rigidity argument excludes them;
their absence is not an input to the double.

Define on $\mathcal U$
\begin{equation}\label{q4:cla:conformal-metrics}
 k_\pm=\left(\frac{1\pm\lambda}{2}\right)^2h_b.
\end{equation}
The four-dimensional scalar conformal formula is
\begin{equation}\label{q4:cla:scalar-transform}
 \Scal_{f^2h}=f^{-3}(-6\Delta_h f+\Scal_h f).
\end{equation}
Thus both $k_+$ and $k_-$ are scalar flat.  In the variables of
\eqref{q4:cla:reduced-variables}, the useful exact identities are
\begin{equation}\label{q4:cla:hyperbolic-factors}
 k_+=\cosh^2(U/2)\gamma,
 \qquad k_-=\sinh^2(U/2)\gamma.
\end{equation}

\begin{lemma}[A zero-mass complete space]\label{q4:cla:double}
Glue a plus and a minus copy along their common attached boundary $S$,
and add one point at the minus asymptotic end.
The resulting space $W$ is a connected orientable smooth
boundaryless four-manifold with a complete $C^2$ scalar-flat metric
$k_0$.  This metric is smooth except possibly at the compact gluing
hypersurface and the added point, and has a distinguished
end satisfying
\begin{equation}\label{q4:cla:zero-end}
 k_0-\delta=O_k(r^{-3})
 \quad\hbox{for every fixed }k.
\end{equation}
No assertion about the other ends of $W$ is needed here.
\end{lemma}

\begin{proof}
The smooth manifold structure is fixed before the metric is smoothed.
At the seam use product charts from the regular base collar, with a
signed normal coordinate on the double.  In the positive boundary-lapse
case this coordinate is signed $\lambda$; in the zero boundary-lapse
case one can use signed base normal distance.  Tangential chart changes
come from $S$, so these charts are smoothly compatible across the seam.
Their overlaps with the original open base are smooth away from the
seam.  The added-point chart below is also smoothly compatible on every
punctured overlap.  Smoothness of this atlas does not assert smoothness
of the reflected metric.

By \eqref{q4:cla:hyperbolic-factors} and
Lemma~\ref{q4:cla:asymptotics}, the plus metric satisfies
$k_+-\delta=O_k(r^{-3})$; in particular its energy is zero.

For the minus end use inversion $x=y/|y|^2$ and put $s=|y|$.
Equation~\eqref{q4:cla:U-expansion} gives
\begin{equation}\label{q4:cla:inverted-lapse}
 \frac{U(x(y))}{s^2}
 =-M_1+H_1(y)+H_2(y)+O_k(s^{2+\epsilon}).
\end{equation}
Differentiating the transformed remainder costs the corresponding
powers of $s^{-1}$.  The Jacobian of inversion is $s^{-2}A(y)$,
where $A(y)=\Id-2yy^t/|y|^2$ is orthogonal.  Consequently the
compactifying expression for $k_-$ is
\begin{equation}\label{q4:cla:inverted-metric}
 \left(\frac{\sinh(U/2)}{s^2}\right)^2
 \left[\Id+A(y)\bigl(\gamma(x(y))-\Id\bigr)A(y)\right].
\end{equation}
The first factor is the product
\[
 \frac14\left(\frac{U}{s^2}\right)^2
       \left(\frac{\sinh(U/2)}{U/2}\right)^2.
\]
It has a $C^2$ extension with positive value $M_1^2/4$ at the origin.
For the second factor, \eqref{q4:cla:gamma-decay} gives
$\gamma(x(y))-\Id=O_k(s^{2+\epsilon})$.
Each derivative of the angular matrix $A$ costs at most one power
of $s^{-1}$, so its conjugated error also has continuous derivatives
through order two vanishing at the origin.  This proves a
nondegenerate $C^2$ extension of \eqref{q4:cla:inverted-metric} across
the added point.  Scalar flatness extends there by continuity.

Proposition~\ref{q4:sta:complete-base} supplies the
even $C^2$ reflection of $h_b$ and the odd $C^2$ reflection of
$\lambda$.  The two conformal factors in
\eqref{q4:cla:conformal-metrics} are one expression
$((1+\lambda)/2)^2$ in that signed lapse.  They therefore glue to a
$C^2$ scalar-flat metric.  Orient the second copy oppositely before
gluing.  The resulting manifold is orientable, and the punctured-ball
chart above extends its orientation over the added point.

It remains to check completeness, including possible escapes toward
interior zeros of $N$.  The plus length factor is at least $1/2$.
On the minus copy its length factor is bounded below away from the
chosen end: if a sequence there had $\lambda\to1$, compactness in
the original exterior would give either a positive-lapse limiting
point with value one, contrary to the strict maximum principle, or
a limiting boundary of the positive component, where $\lambda\to0$.
The latter possibility includes all interior zero loci and the
attachment to $S$.  Thus the complete base ends remain complete for
both conformal metrics.  The seam and the added point are regular
metric neighborhoods.  These observations exhaust the possible
finite-distance escapes by Proposition~\ref{q4:sta:complete-base} and
the completeness of the original exterior.
\end{proof}

\subsection{Zero-mass rigidity without assumptions on the other ends}

We use only the nonnegativity part of a smooth positive-mass theorem.
The precise statement needed is
Theorem~1.2 of Lesourd--Unger--Yau, in arXiv version~1~\cite{LesourdUngerYau2021}:
a complete smooth orientable manifold of dimension $3\le n\le7$,
with nonnegative scalar curvature and a distinguished asymptotically
Schwarzschild end, has nonnegative mass at that end; its other ends
are unrestricted.  Definition~1.9 in that reference requires the
remainder from the Schwarzschild metric to have weighted
$C^{2,\alpha}$ order $n-1$.  There is no spin hypothesis.
The following argument reduces our $C^2$ zero-mass situation to
exactly that smooth nonnegativity statement.

\begin{lemma}[Zero-mass rigidity for the constructed space]
\label{q4:cla:zero-mass}
Let $(W,k_0)$ have the properties in Lemma~\ref{q4:cla:double}.
Then $(W,k_0)$ is isometric to Euclidean $\R^4$.
The isometry is smooth on every region where $k_0$ is smooth,
including each smooth one-sided structure at the seam.
\end{lemma}

\begin{proof}
Suppose first that $\Ric_{k_0}$ is nonzero somewhere.
By continuity it is nonzero at a point in the smooth part of the
metric.  Choose a smooth compactly supported symmetric tensor $p$
in that part for which
\begin{equation}\label{q4:cla:positive-scalar-variation}
 I:=\int_W \Scal'_{k_0}(p)\,dV_{k_0}
   =-\int_W\langle\Ric_{k_0},p\rangle\,dV_{k_0}>0.
\end{equation}
For example, a negative cutoff multiple of $\Ric_{k_0}$ works.
In this proof only, $s>0$ denotes a small metric-variation parameter.
Smooth the compact nonsmooth loci to obtain smooth metrics $\ell_s$
equal to $k_0$ outside one fixed compact set, with
\[
 \ell_s=k_0+s p+o(s)\quad\hbox{in }C^2
\]
on that compact set.  More explicitly, choose fixed nested relatively
compact neighborhoods of the attachment loci, and a smooth cutoff
$\chi$ equal to one on the smaller neighborhood and supported in the
larger one.  Finite-chart mollification and a partition of unity give
a smooth tensor $t_s$ approximating $k_0+s p$ in $C^2$ on the larger
neighborhood with error at most $s^2$.  Set
$\ell_s=\chi t_s+(1-\chi)(k_0+s p)$ there and retain $k_0+s p$
elsewhere.  Where $\chi$ is not one, the old metric is smooth, so this
defines a globally smooth metric.  The perturbation and all smoothing
are confined to a fixed compact set; every other end is unchanged.
For small $s$ the metrics are uniformly comparable to $k_0$ and
therefore complete.  Their scalar curvatures
$R_s=\Scal_{\ell_s}$ have support in a fixed compact set $K$,
satisfy $\|R_s\|_\infty=O(s)$, and obey
\begin{equation}\label{q4:cla:scalar-integral}
 \int_W R_s\,dV_{\ell_s}=sI+o(s)>0.
\end{equation}

We construct a positive smooth conformal factor satisfying
\begin{align}
 -6\Delta_{\ell_s}f_s+R_sf_s&=0,
       & f_s&=1+O(s)\quad\hbox{uniformly on }W,
       \label{q4:cla:conformal-equation}\\
 f_s-1&=O_k(r^{-2})
       &&\hbox{on the distinguished end},
       \label{q4:cla:factor-decay}\\
 \int_{S_R}\partial_{\nu_{\ell_s}}f_s\,dA_{\ell_s}
       &=\frac16\int_W R_sf_s\,dV_{\ell_s}
       &&\hbox{for every sufficiently large }R.
       \label{q4:cla:single-end-flux}
\end{align}
In particular, the flux in \eqref{q4:cla:single-end-flux} is not being
obtained by assuming zero flux at unspecified ends.

Take connected smooth compact exhaustions $D_j$ whose boundary in
the distinguished end is a large coordinate sphere $S_{R_j}$.
The remaining boundary components lie outside each prescribed
compact subset for sufficiently large $j$.  Such an exhaustion is
obtained by exhausting the complement of a fixed end neighborhood,
joining to full coordinate annuli, and smoothing the compact joins.
Impose $f=1$ on $S_{R_j}$ and homogeneous Neumann conditions on every
other boundary component.  These boundary components are disjoint,
so the mixed problem has no interface corners.

We first record a uniform estimate for the compactly supported
potential.  If $v$ has zero Dirichlet value on $S_{R_j}$, then for
every fixed compact set $K'$ in a fixed connected core neighborhood,
\begin{equation}\label{q4:cla:anchored-estimate}
 \|v\|_{L^2(K',\ell_s)}
       \le C_{K'}\|\nabla v\|_{L^2(D_j,\ell_s)},
\end{equation}
with a constant independent of sufficiently large $j$ and small $s$.
To see this first on a fixed end annulus, integrate along coordinate
rays to the Dirichlet sphere.  In Euclidean polar coordinates,
\[
 |v(t,\theta)|^2
 \le\left(\int_t^{R_j}r^3|\partial_rv(r,\theta)|^2\,dr\right)
      \left(\int_t^{R_j}r^{-3}\,dr\right),
 \qquad \int_t^{R_j}r^{-3}\,dr\le\frac1{2t^2}.
\]
Integration over the unit sphere and over the fixed annulus controls
its $L^2$ norm by the full gradient energy.  On a fixed connected
neighborhood joining that annulus to $K'$, the Poincare inequality
with the annulus as an anchor propagates the estimate.  Uniform
metric comparison on that neighborhood and the tail makes the
constants uniform in $s$.

For $v=f-1$ the variational problem is
\begin{equation}\label{q4:cla:mixed-weak}
 \int_{D_j}\left(\langle\nabla v,\nabla w\rangle
                   +\frac{R_s}{6}vw\right)dV_{\ell_s}
       =-\frac16\int_{D_j}R_sw\,dV_{\ell_s},
\end{equation}
for tests $w$ vanishing on the Dirichlet sphere.  By
\eqref{q4:cla:anchored-estimate} with a fixed neighborhood of $K$,
the left quadratic form is bounded below by
$(1-Cs)\|\nabla v\|_2^2$, and the right side is bounded in absolute
value by $Cs\|\nabla w\|_2$.  On each finite domain the gradient norm
is a Hilbert norm on the Dirichlet test space.  Lax--Milgram therefore
solves \eqref{q4:cla:mixed-weak} for small $s$, with
\begin{equation}\label{q4:cla:mixed-energy}
 \|\nabla v\|_{L^2(D_j)}=O(s),\qquad
 \|v\|_{L^2(K')}=O(s)
\end{equation}
for every fixed compact $K'$ as above.  Smooth elliptic regularity
holds for the finite-domain solutions.

These estimates also give a uniform supremum bound near $K$.
First, interior $W^{2,2}$ estimates for
$\Delta_{\ell_s}v=(R_s/6)(1+v)$ give $O(s)$ on a slightly smaller
fixed neighborhood.  In dimension four this implies every finite
local $L^p$ bound, though not yet an $L^\infty$ bound.
Applying the equation once more gives $W^{2,p}=O(s)$ for $p>2$,
and hence $\|v\|_\infty=O(s)$ there.
The constants are uniform because the metrics have uniform $C^2$
bounds and uniform ellipticity on those neighborhoods, and
$R_s=O(s)$ in $L^\infty$.  Derivative bounds on the smoothing error
beyond those needed here are unnecessary.

Outside a fixed neighborhood of $K$, the function $v$ is harmonic.
Its values at the inner boundary are bounded by $Cs$, its outer
Dirichlet values are zero, and its other outer normal derivatives
are zero.  Here is a global weak test that avoids assumptions on the
remaining components.  Take $(v-Cs)_+$, which vanishes near $K$, and
extend it by zero through that neighborhood.  It is an admissible test
in \eqref{q4:cla:mixed-weak}; both terms containing $R_s$ vanish, leaving
the integral of its squared gradient equal to zero.  Connectedness of
$D_j$ and its zero Dirichlet trace imply $(v-Cs)_+=0$ everywhere.
The negative excess gives the other bound.  Thus $|v|\le Cs$
throughout $D_j$, with no geometry of the other ends used.
On the distinguished end, a multiple of
$r^{-2}(1-r^{-\eta})$, with $0<\eta<1$, is a positive strictly
superharmonic barrier sufficiently far out for the metric
\eqref{q4:cla:zero-end}.  Comparison on the truncated annuli gives
the uniform bound $|v|\le C_s r^{-2}$ there.

For each fixed small $s$, pass to a subsequence converging smoothly
on compact subsets as $j\to\infty$.  The limit yields
\eqref{q4:cla:conformal-equation} and the zeroth-order decay in
\eqref{q4:cla:factor-decay}.  Rescaled end estimates and differentiation
give all its fixed differentiated orders.  Positivity follows from
the uniform bound $f_s=1+O(s)$.
For a fixed sphere $S_R$ outside the support of $R_s$, integrate
the finite-domain equation on
\[
 B_{j,R}=D_j\setminus\{x\text{ in the distinguished end}:r>R\}.
\]
Its boundary consists of $S_R$ and only artificial boundaries with
homogeneous Neumann data; the Dirichlet sphere $S_{R_j}$ has been
removed.  For large $j$ it contains the whole support of $R_s$, so
\[
 \int_{S_R}\partial_{\nu_{\ell_s}}f_{s,j}\,dA_{\ell_s}
       =\frac16\int_{D_j}R_sf_{s,j}\,dV_{\ell_s}.
\]
The curvature has fixed compact support.  Local convergence thus
passes the derivative on the fixed sphere and the integral on that
fixed support to the limit, proving \eqref{q4:cla:single-end-flux} exactly.
No integration over a limiting noncompact inner region is involved.
The limiting factor may have harmonic behavior at the other ends;
neither its limiting values nor individual end fluxes are required.

The metric $f_s^2\ell_s$ is smooth, complete and scalar flat by
\eqref{q4:cla:scalar-transform} and the uniform positive bounds for
$f_s$.  Its distinguished end is asymptotically Schwarzschild in
the precise sense required by the stated positive-mass theorem.
Indeed, on that end $\ell_s=k_0=\delta+O_k(r^{-3})$ and
$\Delta_{\ell_s}f_s=0$.  Combining this with
\eqref{q4:cla:factor-decay} gives
$\Delta_\delta(f_s-1)=O_k(r^{-7})$.
The multipole estimate \eqref{q4:cla:multipole} therefore implies
\begin{equation}\label{q4:cla:factor-monopole}
 f_s=1+\frac{a_s}{r^2}+O_k(r^{-3}),\qquad
 f_s^2\ell_s=(1+a_s/r^2)^2\delta+O_k(r^{-3}).
\end{equation}
The differentiated bounds supply the weighted $C^{2,\alpha}$
remainder of order three, required in dimension four.

On the other hand, \eqref{q4:cla:scalar-integral} and the uniform
bound in \eqref{q4:cla:conformal-equation} give
\[
 \int_W R_sf_s\,dV_{\ell_s}=sI+o(s)>0.
\]
Thus \eqref{q4:cla:single-end-flux} has strictly positive outward
flux.  Physical and Euclidean fluxes have the same limit by the
falloff.  From \eqref{q4:cla:factor-monopole}, that limit is
$-2\omega_3a_s$, so $a_s<0$.  Equivalently, direct computation with
the prescribed energy normalization gives
\begin{equation}\label{q4:cla:negative-mass}
 E(f_s^2\ell_s)
   =-\frac1{\omega_3}
        \lim_{R\to\infty}\int_{S_R}\partial_r f_s\,dA_\delta
   =2a_s<0.
\end{equation}
The background end contributes zero energy, and products of its
error with the conformal error contribute zero limiting flux.
This contradicts the smooth arbitrary-ends positive-mass theorem:
the dimension is four, the manifold is orientable and boundaryless,
the metric is complete and smooth, its scalar curvature vanishes,
and \eqref{q4:cla:factor-monopole} has exactly the required distinguished
end.  Hence $\Ric_{k_0}=0$.

For clarity concerning regularity, a $C^2$ Ricci-flat metric becomes
smooth in harmonic coordinates by the elliptic Ricci equation;
this is the harmonic-coordinate regularity statement of
DeTurck--Kazdan~\cite[Theorem~5.2]{DeTurckKazdan1981} for $C^2$ Einstein
metrics in dimension at least three.
We may consequently apply Bishop--Gromov comparison and its
rigidity statement~\cite{Petersen2016}.  The asymptotic volume ratio
is at least the Euclidean value using the distinguished end alone:
paths from a fixed base point to a fixed end sphere, followed by
coordinate rays, have length $r+o(r)$ uniformly in angle, and the
end volume form tends to its Euclidean value.  Thus large metric
balls contain asymptotically Euclidean coordinate annuli of the
corresponding radius.  Ricci nonnegativity gives the opposite
inequality for the volume ratio.  Its value is therefore one, and
the equality case of Bishop--Gromov implies that $W$ is Euclidean
space globally.

The isometry is smooth wherever the original metric is smooth.
Its smoothness in each one-sided seam structure can also be seen
directly, without asserting that the reflected metric was smooth.
The $C^2$ metric has $C^1$ connection coefficients in the original
seam charts.  A flat parallel coframe extends across the seam by
parallel transport, with at least $C^1$ dependence on its base point.
In coordinates each member $\alpha$ satisfies
\[
 \partial_j\alpha_i=\Gamma^k_{ji}\alpha_k.
\]
On each closed one-sided chart the connection coefficients are smooth
up to the boundary.  This identity inductively makes the coframe
smooth up to that boundary.  Its closed forms integrate to Euclidean
coordinates agreeing with the isometry on the open side after a fixed
Euclidean motion, hence also at the seam by continuity.  This proves
the asserted smoothness separately in both one-sided structures.
\end{proof}

The complete double is now Euclidean. This is stronger than a local
static classification: it also excludes any complete base end hidden at
an interior zero of the original Killing norm. We use that conclusion
first, and only then identify the boundary sphere and the explicit lapse.

\subsection{The original positive component and its spherical boundary}

\begin{proposition}[Global static classification]
\label{q4:cla:static-classification}
The positive component is the whole original open exterior:
$\mathcal U=\operatorname{int}\Omega$ and $N>0$ there.
Its attached boundary is $S$.  There is a global isometry of the
static base onto the Schwarzschild--Tangherlini spatial exterior of
mass $m$, under which
\begin{equation}\label{q4:cla:ST-base}
 h_b=\left(1+\frac{m}{2\rho^2}\right)^2
          (\dd\rho^2+\rho^2 g_{\mathbb S^3}),\qquad
 \lambda=\frac{1-m/(2\rho^2)}{1+m/(2\rho^2)},
 \qquad \rho\ge\sqrt{m/2}.
\end{equation}
The isometry is smooth on the open base and in its one-sided regular
boundary structure.  The attachment is homeomorphic to the original
attachment of $S$ in $\Omega$ and restricts smoothly on $S$.
In particular the open base is simply connected and $S$ is a round
three-sphere in its induced metric.  The base reflection is smooth
in signed $\lambda$ in the positive boundary-lapse case; its angular
coordinates can be taken constant along base normal geodesics.
\end{proposition}

\begin{proof}
By Lemmas~\ref{q4:cla:double} and~\ref{q4:cla:zero-mass}, the constructed
space $W$ is Euclidean.  Suppose a sequence in $\mathcal U$ approaches
an interior zero of $N$ in the original compact part of $\Omega$.
It escapes every compact subset of $W$: the zero locus is not a
point of the plus copy, it cannot converge to $S$, and the only
new point lies at the minus end in a neighborhood disjoint from
that sequence.  Yet the sequence remains outside a fixed distant
tail of the distinguished end of $W$.

This is impossible in Euclidean space.  A coordinate cross-sphere
in the distinguished end is a compact embedded sphere.  Its tail
has only that sphere as its frontier and is the unbounded component
of its complement.  Jordan--Brouwer separation makes the other
side bounded; its closure is compact.  Thus the complement of the
tail is compact.  This use of separation does not require a smooth
Schoenflies theorem in dimension four.

It follows that $\mathcal U$ has no boundary in
$\operatorname{int}\Omega$.  The latter is connected, so
$\mathcal U=\operatorname{int}\Omega$.  We now identify the plus
and minus copies in the Euclidean double explicitly.

At $S$, under the change $k_+=f^2h_b$ with
$f=(1+\lambda)/2$, the shape operator transforms by
\[
 \mathcal S_{k_+}
       =f^{-1}\bigl(\mathcal S_{h_b}
                   +\nu_{h_b}(\log f)\Id\bigr).
\]
Using \eqref{q4:cla:boundary-data}, its value is $2\kappa\Id$ for
the normal toward the plus side.  In Euclidean coordinates write
$z$ for position and $\nu$ for that normal.  With
$d_*=(2\kappa)^{-1}$, the tangential derivative of $z-d_*\nu$
vanishes.  Connectedness of $S$ makes this vector a constant center.
The image of $S$ lies on the sphere of radius $d_*$ about that
center, and is both open there by local immersion and closed by
compactness.  It is the entire round sphere.  Embeddedness follows
from the global ambient isometry.  The plus side, which contains
the distinguished infinity, is its Euclidean exterior.

Set $\psi=2/(1+\lambda)$.  Then $h_b=\psi^2k_+$.
Scalar flatness and \eqref{q4:cla:scalar-transform} show that $\psi$
is harmonic on that Euclidean exterior.  It is equal to two on
the sphere and tends to one at infinity.  Uniqueness by the
maximum principle therefore gives, in centered Euclidean radius
$\rho$,
\begin{equation}\label{q4:cla:psi-and-lapse}
 \psi=1+\frac{d_*^2}{\rho^2},\qquad
 \lambda=\frac{1-d_*^2/\rho^2}{1+d_*^2/\rho^2}.
\end{equation}
Since $\kappa=1/\sqrt{2m}$, we have $d_*^2=m/2$,
which proves \eqref{q4:cla:ST-base}.  For an invariant check on the
coefficient from Lemma~\ref{q4:cla:asymptotics}, the flux of the
$\gamma$-harmonic function $U$ is independent of the end cross-section.
Its limit is $2\omega_3M_1$ in the harmonic coordinates and
$2\omega_3m$ in the explicit coordinates of \eqref{q4:cla:psi-and-lapse}.
The divergence theorem between homologous cross-sections therefore
gives $M_1=m$, without identifying the two asymptotic charts.
The area radius
\[
 r=\rho\left(1+\frac{m}{2\rho^2}\right)
\]
satisfies
\begin{equation}\label{q4:cla:area-radius-base}
 N=1-\frac{2m}{r^2},\qquad
 h_b=N^{-1}\dd r^2+r^2g_{\mathbb S^3}
       \quad\hbox{on }\operatorname{int}\Omega.
\end{equation}
The apparent degeneracy of the area-radius expression at the
boundary is removed by the isotropic radius or the regular base
collar.  Its boundary radius is $r_0=\sqrt{2m}$.

The smooth one-sided boundary assertion follows from
Lemma~\ref{q4:cla:zero-mass} and the regular base collar of
Proposition~\ref{q4:sta:complete-base}.  In the positive boundary-lapse
case this collar is expressed in $\lambda=\sqrt N$, with a smooth
reflection; replacing an original defining coordinate $N$ by
$\sqrt N$ changes the smooth boundary coordinate but not the
underlying topology or the smooth structure on $S$ itself.
In the zero boundary-lapse case the base collar is already smooth
in the original one-sided structure.  The spherical angular
coordinates in \eqref{q4:cla:ST-base} are constant along the base
normal geodesics from $S$.  In a reflected collar their normal
projection to $S$ is smooth and invariant under reflection.
These facts give precisely the regular boundary identification
needed to recover the hypersurface in horizon coordinates.
Finally the Euclidean exterior of a ball in dimension four is
simply connected, proving the topological assertions.
\end{proof}
